% Nouveau EGA V (cote 25): archive cote 25, site page 158 (PDF page 159).
% English translation of the checked French annotated master.
% AI-translated by GPT-6 Astra (OpenAI, ChatGPT Pro mode).
\scanpage{25/158}
\var{\margin{new EGA \uline{V}}}{\placedmargin{new EGA \uline{V}}}{marginal note placed}{ega25_25_158-v0001} \mspage{\hand{1}}

IV 17.15. \uline{Smooth forms}; \var{\inserted{augmentations and elementary singularities.}}{augmentations and elementary singularities.}{insertion}{ega25_25_158-v0002} \var{\deleted{smooth quadratic forms}}{}{deletion}{ega25_25_158-v0003}

NB I have just realized \hand{\textsuperscript{\unsure{$\dagger$}}}that the terminology introduced in
my account \var{\inserted[typed]{SUP}}{SUP}{insertion}{ega25_25_158-v0005} of Supersingularity is unreasonable, and conflicts
in particular with the accepted terminology. (In any case, you \unsure{must} have
noticed that in def.1 of this text one must read ``degenerate'' instead of
``nondegenerate''). The absurdity comes from varieties of even dimension
in char.2, \var{\xout{such a variety}}{}{typewriter strike-out}{ega25_25_158-v0007} defined, for example, by an equation $\hand{\phi}=0$
in an ambient variety of odd dim, for with the terminology of my
notes, such a variety cannot have an \var{\xout{sing zero}}{}{typewriter strike-out}{ega25_25_158-v0008} \hand{``}ordinary
singular point\hand{''} (= ordinary quadratic singularity), i.e. $\hand{\phi}$ cannot have
an \hand{``}ordinary singular zero\hand{''}, owing to the fact that in char 2, a quadratic form
in an odd number of variables is always ``degenerate''. Yet everyone
has always considered that even in char. 2, the origin of the affine cone
$x^2 + yz = 0$ is an ordinary quadratic singularity. In the present
notes, I bring out the notion of a smooth quadratic \unsure{f|o}\unsure{o|f}rm (or an ``ordinary'' one),
(on a loc free Module of finite type), in such a way that
nondegenerate \hand{$\Longrightarrow$} smooth, the converse being true if E has even rank
\uline{or} the residue char of the base S are all $\hand{\neq} 2$ . Having
introduced this notion, I propose to keep the terminology ``supersingular'' (which does not
conflict with any accepted terminology) \var{\xout{as in my text}}{}{typewriter strike-out}{ega25_25_158-v0010} of SUP,
as corresponding to the notion of a nondegenerate quadratic form; we
shall likewise speak of a supersingular point of a loc noeth. prescheme, or a
geometrically supersingular point for a prescheme loc of finite type over a
field, in the same sense (by considering the hom $\mathrm{Sym}(\underline{\unsure{M|m}}/\underline{m}^2) \hand{\longrightarrow} \mathrm{gr}_{\hand{\underline{\textcolor{black}{m}}}}(A)$ and
\unsure{by} saying that the kernel cannot be generated by a nondegenerate quadratic
form); thus what was wrongly called an ``ordinary singular
zero'' in SUP will be called a ``singular nonsupersingular zero'', or a ``nondegenerate quadratic

\mspage{\hand{158}}
\tnote{25-158-1}{The source abbreviations def., char., dim, loc, noeth. and hom mean definition, characteristic, dimension, locally, noetherian and homomorphism. SUP is retained as the source's designation of its account of supersingularity. The colloquial canular is rendered absurdity. The doubt in must comes from the final s of French as; the two doubts in form retain the source's f/o and o/f candidates, and by renders the doubtful en/ex. The handwritten accents and plural or final letters have no separate English spelling effect.}
\tnote{25-158-2}{The retouched t in noter, d in ordinaire, and f-to-r correction in dégénérée do not change the English words and receive no English variants; the old \unsure{d} remains doubtful. The deleted isolated x in the heading has no English lexical content. The final quoted term continues on the next scan; its adjectives precede singular zero in English. The source's stated correspondences and negations are retained.}
